% ============================================================
\begin{document}
\selectlanguage{english}

% ============================================================
%  FRONT MATTER
% ============================================================
\begin{titlepage}
  \centering

  {\Large\bfseries Technical University of Crete\\[0.3em]
  School of Electrical and Computer Engineering}\\[0.5em]
  {\large Diploma Thesis}\\[0.8cm]

  \rule{\linewidth}{0.5mm}\\[0.4cm]
  {\LARGE\bfseries Tracking the Lifetime of Malicious Indicators\\[0.3em]
  in Public Threat Intelligence Feeds}\\[0.4cm]
  \rule{\linewidth}{0.5mm}\\[0.5cm]

  {\large\itshape
  Παρακολούθηση της διάρκειας ζωής κακόβουλων δεικτών\\[0.2em]
  σε ροές πληροφοριών για δημόσιες απειλές}\\[0.9cm]

  \begin{tabular}{@{}l@{\hspace{3cm}}l@{}}
    \textit{Author:} & \textit{Thesis Committee:} \\
    Dimitris Keramidas & [Supervisor Name] (Supervisor) \\
     & [Committee Member 2] \\
     & [Committee Member 3] \\
  \end{tabular}\\[0.8cm]

  \includegraphics[width=2cm]{tuc_logo.png}\\[0.8cm]

  {\itshape A thesis submitted in fulfillment of the requirements\\
  for the diploma of Electrical and Computer Engineer\\[0.3em]
  in the\\[0.3em]
  School of Electrical and Computer Engineering}\\[0.6cm]

  \vfill
  {\large Chania, Crete, Greece \\ \the\year}
\end{titlepage}

\cleardoublepage

% --- Copyright / Declaration ---
\thispagestyle{empty}
\vspace*{\fill}
\begin{center}
  \textcopyright\ \the\year\ Dimitris Keramidas. All rights reserved.
\end{center}
\vspace*{\fill}
\cleardoublepage

% --- Abstract (English) ---
\chapter*{Abstract}
\addcontentsline{toc}{chapter}{Abstract}

Public threat intelligence feeds are a primary resource for identifying malicious activity on the
internet, yet little is known about how they respond over time to malicious servers found by other
means. This thesis tracks IP addresses classified as malicious by TLScope --- an active TLS
fingerprinting system that identifies malicious servers without inspecting payload content ---
across VirusTotal (and the 90--100 engines it aggregates), AbuseIPDB, ThreatFox, Censys, Cinsscore and
OpenPhish, sampling every IP every two to three days for 33--39 days after TLScope's verdict.

Across seven collection waves covering 15\,837 IPs, \textbf{82.5\% were never flagged by any
provider}. Of the 2\,770 that were, 47.7\% were already flagged on the first sample and 41.2\% were
flagged later, typically after one to three weeks (median 15 days). Once raised, most flags were
still in place when observation ended (64.6\% of all engine--IP detections), so measured indicator
lifetimes are lower bounds; the one engine with broad coverage, Criminal IP, changed behaviour
mid-study, from holding its flags for the whole window to dropping them within one to two weeks.
Agreement is sparse: 83.1\% of detected IPs are flagged by exactly one engine, although engines
co-detect more often than independent chance would predict, so the low overlap reflects how few IPs
each engine flags rather than engines covering disjoint niches. VirusTotal's engines account for
95.5\% of all detections. A data-quality audit of the collection itself is reported alongside the
results, including a provider deactivated mid-study and a pilot wave excluded for a classification
bug.

\bigskip\noindent
\textbf{Keywords:} threat intelligence, indicator of compromise, malicious IP, VirusTotal,
TLS fingerprinting, detection lag, indicator lifetime, measurement study

\cleardoublepage

% --- Abstract (Greek) ---
\selectlanguage{greek}
\chapter*{Περίληψη}
\addcontentsline{toc}{chapter}{Περίληψη}

Οι δημόσιες ροές πληροφοριών για απειλές αποτελούν βασικό εργαλείο για τον εντοπισμό κακόβουλης
δραστηριότητας στο διαδίκτυο, ωστόσο ελάχιστα είναι γνωστά για το πώς αντιδρούν με την πάροδο του
χρόνου σε κακόβουλους διακομιστές που έχουν εντοπιστεί με άλλα μέσα. Η παρούσα εργασία
παρακολουθεί διευθύνσεις \foreignlanguage{english}{IP} που έχουν ταξινομηθεί ως κακόβουλες από το
\foreignlanguage{english}{TLScope}, ένα σύστημα ενεργού αποτυπώματος \foreignlanguage{english}{TLS},
στους παρόχους \foreignlanguage{english}{VirusTotal} (και τις 90--100 μηχανές που συγκεντρώνει),
\foreignlanguage{english}{AbuseIPDB}, \foreignlanguage{english}{ThreatFox},
\foreignlanguage{english}{Censys}, \foreignlanguage{english}{Cinsscore} και
\foreignlanguage{english}{OpenPhish}, με δειγματοληψία κάθε δύο έως τρεις ημέρες για 33--39 ημέρες
μετά την ταξινόμηση.

Σε επτά κύκλους συλλογής με 15.837 διευθύνσεις, το \textbf{82,5\% δεν επισημάνθηκε ποτέ από
κανέναν πάροχο}. Από τις 2.770 που επισημάνθηκαν, το 47,7\% ήταν ήδη επισημασμένο στο πρώτο δείγμα
και το 41,2\% επισημάνθηκε αργότερα, συνήθως μετά από μία έως τρεις εβδομάδες (διάμεσος 15 ημέρες).
Οι περισσότερες επισημάνσεις παρέμεναν ενεργές στο τέλος της παρακολούθησης (64,6\%), επομένως οι
μετρούμενες διάρκειες ζωής αποτελούν κάτω φράγματα. Η συμφωνία μεταξύ μηχανών είναι αραιή: το
83,1\% των διευθύνσεων επισημαίνεται από μία μόνο μηχανή, αν και οι μηχανές συμπίπτουν συχνότερα από
ό,τι θα προέβλεπε η τυχαία ανεξαρτησία, επομένως η χαμηλή επικάλυψη οφείλεται στο ότι κάθε μηχανή
επισημαίνει ελάχιστες διευθύνσεις. Οι μηχανές του \foreignlanguage{english}{VirusTotal} ευθύνονται
για το 95,5\% των ανιχνεύσεων. Η εργασία περιλαμβάνει επίσης έλεγχο ποιότητας της ίδιας της
συλλογής δεδομένων.

\bigskip\noindent
\textbf{Λέξεις-κλειδιά:} πληροφορίες απειλών, δείκτης παραβίασης, κακόβουλη
\foreignlanguage{english}{IP},
\foreignlanguage{english}{VirusTotal}, αποτύπωμα \foreignlanguage{english}{TLS}, καθυστέρηση
ανίχνευσης, διάρκεια ζωής δείκτη, μελέτη μετρήσεων

\selectlanguage{english}
\cleardoublepage

% --- Acknowledgements ---
\chapter*{Acknowledgements}
\addcontentsline{toc}{chapter}{Acknowledgements}

I would like to thank my supervisor for their guidance throughout this project,
and the authors of TLScope for providing the classified IP datasets that form the
empirical foundation of this work.

\cleardoublepage

% --- Table of Contents ---
\tableofcontents
\cleardoublepage

\listoffigures
\addcontentsline{toc}{chapter}{List of Figures}
\cleardoublepage

\listoftables
\addcontentsline{toc}{chapter}{List of Tables}
\cleardoublepage

\glsaddall
\printglossary[type=\acronymtype, title={List of Abbreviations}, nonumberlist]
\cleardoublepage

% ============================================================
%  MAIN MATTER
% ============================================================

% ============================================================
\chapter{Introduction}
\label{ch:introduction}
% ============================================================

\section{Motivation}

Public threat intelligence feeds --- continuously updated lists of known-malicious IP addresses,
domains and URLs --- sit underneath much of operational network security. Security operations
centres, firewalls, intrusion detection systems and endpoint products consume them to block or flag
traffic associated with malware, botnet command-and-control (C2) infrastructure, phishing and abuse.
The practical value of a feed rests on two properties: its \textbf{coverage}, the fraction of
genuinely malicious indicators it lists, and its \textbf{timeliness}, how quickly it lists them and
how long it keeps them listed.

Coverage has been measured before, usually by comparing feeds against one another
(Chapter~\ref{ch:related_work}). Timeliness is harder, because measuring it requires an independent
and \emph{dated} source of truth: one needs to know when a server was known to be malicious in order
to say how long a feed took to notice. This thesis obtains that from TLScope, an active TLS
fingerprinting system that classifies servers as malicious from their handshake behaviour alone and
records when it did so. Each TLScope batch therefore gives a set of servers known to be malicious on
a known date, which can then be followed forward through the public feeds.

\section{Research Questions}

\begin{enumerate}[label=\textbf{RQ\arabic*.}]
  \item \textbf{Detection lag:} After TLScope identifies a server as malicious, how long does it
    take for public threat intelligence providers to flag its IP address --- if they do at all?
  \item \textbf{Indicator persistence:} Once a provider flags an IP, how long does the flag
    remain in place?
  \item \textbf{Inter-provider agreement:} Do different providers and engines flag the same IPs,
    and is the observed overlap more or less than their individual detection rates alone would
    produce?
\end{enumerate}

\section{Approach}

Seven waves of TLScope-classified malicious IPs (15\,837 IPs, 1\,596--2\,728 per wave) were tracked
between December 2025 and September 2026. Every IP was queried against six providers every three
days (two days in one wave) for 33--39 days, starting on the day TLScope produced its verdict:

\begin{itemize}
  \item \textbf{VirusTotal} --- an aggregator reporting the verdicts of 90--100 third-party
    blocklist and security engines;
  \item \textbf{AbuseIPDB} --- community-reported abuse confidence score;
  \item \textbf{ThreatFox} --- abuse.ch's malware IOC database;
  \item \textbf{Censys} --- internet-wide scanner with C2 labels (queried in the first three and a
    half waves only, \S\ref{sec:provider_availability});
  \item \textbf{Cinsscore} --- the CINS Army daily IP blocklist;
  \item \textbf{OpenPhish} --- phishing URL feed, resolved to IPs by DNS.
\end{itemize}

An eighth, earlier wave was collected while the collection tool was still under development and is
excluded from every result in this thesis (\S\ref{sec:data_quality}).

\section{Contributions}

\begin{enumerate}
  \item \textbf{A longitudinal, ground-truth-anchored dataset}: 15\,837 TLScope-identified
    malicious IPs, each sampled on average 10--17 times across six providers over 33--39 days from the date of
    its TLScope classification.
  \item \textbf{A coverage measurement}: 82.5\% of these IPs are never flagged by any provider
    within the window. A sensitivity check on one wave shows that this figure depends strongly on
    window length (68.2\% when the same wave is followed for 60 days instead of 39), so it is
    reported as a window-specific quantity rather than an absolute one.
  \item \textbf{A lag measurement in days}: detection is bimodal --- an IP is either already
    flagged when TLScope finds it (47.7\% of detected IPs) or flagged one to three weeks later
    (median 15 days) --- and late detections arrive in bulk listing events as often as one by one.
  \item \textbf{A persistence measurement with explicit censoring}: 64.6\% of detections are still
    in place when observation ends, so lifetimes are reported as lower bounds, and a change in one
    major engine's retention behaviour during the study is identified.
  \item \textbf{A corrected agreement analysis}: comparing observed overlap against an independence
    model shows that engines co-detect \emph{more} often than chance (1.46--2.10$\times$ the
    expected multi-engine share in every wave); low agreement is a consequence of sparse individual
    coverage, not of specialisation.
  \item \textbf{A data-quality audit of the collection} documenting a provider deactivated
    mid-study, a classification bug that invalidated the pilot wave, and the incomplete-response
    rate, so that the results above can be weighed against how they were produced.
  \item \textbf{Open-source tooling} (\texttt{NetScan}) for reproducible multi-provider IP
    tracking.
\end{enumerate}

\section{Thesis Structure}

Chapter~\ref{ch:background} provides background on TLS fingerprinting and threat intelligence
feeds. Chapter~\ref{ch:related_work} surveys related work. Chapter~\ref{ch:methodology} describes the
data collection, the measurement definitions and the data-quality audit.
Chapter~\ref{ch:providers} details each provider and Chapter~\ref{ch:implementation} the NetScan
implementation. Chapter~\ref{ch:results} presents the results and Chapter~\ref{ch:discussion}
answers the research questions and sets out the threats to validity.
Chapter~\ref{ch:future_work} outlines future directions and Chapter~\ref{ch:conclusion} concludes.

% ============================================================
\chapter{Background}
\label{ch:background}
% ============================================================

\section{TLS and Encrypted Malicious Infrastructure}

The Transport Layer Security (TLS) protocol is the cryptographic foundation of HTTPS and many
other application protocols. By 2022, more than 85\% of malware used encrypted communication
channels, making payload-based detection infeasible without decrypting
traffic~\cite{theofanous2024tlscope}. TLS metadata --- cipher suite selection, protocol version,
extension handling, certificate properties --- remains visible even when the payload is encrypted,
and provides an alternative detection surface.

\subsection{JARM Fingerprinting}

JARM~\cite{jarm2020} is a widely used active TLS fingerprinting technique. It sends ten specially
crafted \texttt{ClientHello} messages to a server, varying parameters such as cipher suite ordering
and extensions. Each server's TLS stack responds according to its implementation, and the pattern
of responses across the ten probes is hashed into a 62-character fingerprint. Servers running the
same TLS library and configuration produce the same JARM hash, enabling large-scale infrastructure
clustering.

\subsection{TLScope}

TLScope~\cite{theofanous2024tlscope} extends JARM-style fingerprinting with a full machine learning
pipeline. Rather than collapsing responses to a hash, it treats the full per-probe response vector
as a high-dimensional feature input, preserving inter-probe correlations, and adds certificate
metadata, protocol capability flags and server behaviour under adversarial \texttt{ClientHello}
configurations. The resulting classifier reports 91\% precision and 95\% recall for malicious
server detection, and 99\% precision and recall for botnet family identification.

\section{Threat Intelligence Feeds}

Threat intelligence feeds are curated lists of indicators of compromise (IOCs). IP-based feeds
typically assign a reputation score or a binary malicious/benign label to individual addresses.
Feed types relevant to this thesis include:

\begin{description}
  \item[Multi-engine aggregates] Services that query many independent engines and report each
    verdict (VirusTotal). An aggregate is not a feed in its own right; its coverage is the union
    of its engines' coverage.
  \item[Community-reported feeds] Aggregated abuse reports from network operators and automated
    systems (AbuseIPDB).
  \item[Specialised IOC databases] Curated databases of malware-related indicators (ThreatFox).
  \item[Active scanning feeds] Labels derived from continuous internet-wide scanning (Censys).
  \item[Blocklists] Lightweight, regularly regenerated IP lists (Cinsscore).
  \item[Phishing URL feeds] Domain-based phishing indicators, usable as IP indicators only after
    DNS resolution (OpenPhish).
\end{description}

\section{Indicator Lifecycle}
\label{sec:lifecycle}

From the point of view of a single provider, an indicator passes through three phases:
\begin{enumerate}
  \item \textbf{Pre-detection:} the server is malicious but the provider does not yet list it.
  \item \textbf{Active detection:} the provider lists the IP.
  \item \textbf{Expiry:} the provider removes the listing --- because the threat ended, because a
    retention policy aged it out, or because a classifier changed its verdict.
\end{enumerate}
Detection lag is the duration of phase~1 as seen from the moment the server was independently
known to be malicious; indicator lifetime is the duration of phase~2. A finite observation window
cuts both short: an IP not yet detected when observation ends has an unknown lag, and a flag still
in place when observation ends has a lifetime of \emph{at least} the observed span. Both forms of
right-censoring are central to how the results in Chapter~\ref{ch:results} must be read.

% ============================================================
\chapter{Related Work}
\label{ch:related_work}
% ============================================================

\section{TLScope: the Parent Project}

TLScope \cite{theofanous2024tlscope} is not a related work surveyed alongside the papers below ---
it is the \textbf{parent project} this thesis is built directly on top of, and its output is this
thesis's ground truth rather than a point of comparison. TLScope actively initiates TLS connections
to servers and sends exhaustively crafted \texttt{ClientHello} messages that progressively force
each server to select alternative cipher suites, extending JARM-style fingerprinting
(Chapter~\ref{ch:background}) with a full ML pipeline. The system achieves 91\% precision and 95\%
recall for malicious server detection and 99\%/99\% for botnet family identification, and its
output --- dated batches of IP addresses classified as malicious --- is what
Chapter~\ref{ch:methodology} tracks across the public feeds.

The two projects answer different halves of the same underlying question, as
Table~\ref{tab:tlscope_netscan} makes explicit:

\begin{table}[h]
  \centering
  \caption{Division of questions between TLScope and this thesis.}
  \label{tab:tlscope_netscan}
  \begin{tabularx}{\linewidth}{Xl}
    \toprule
    \textbf{Question} & \textbf{Answered by} \\
    \midrule
    Is this server's TLS behaviour indicative of malicious hosting? & TLScope \\
    What botnet family does this server belong to?                 & TLScope \\
    How quickly do public feeds detect what TLScope already found?  & This thesis \\
    How long do providers keep an indicator flagged once detected?  & This thesis \\
    Do providers agree on the same IPs, or cover distinct subsets?  & This thesis \\
    \bottomrule
  \end{tabularx}
\end{table}

Every paper surveyed in the remainder of this chapter is assessed by its relation \emph{to} this
pairing, not treated as a peer of TLScope (a curated version of the survey, cross-referenced against
the Future Work chapter, is kept alongside the source PDFs in \texttt{wiki/papers/notes.md}).

\section{Measurement Studies of Threat Intelligence Feeds}
\label{sec:rw_ti_measurement}

The closest prior work measures threat intelligence feeds directly rather than proposing a new
detector, and it establishes the baseline expectation for RQ3.

Kührer, Rossow and Holz \cite{kuhrer2014paintitblack} analysed 15 public malware blacklists and 4
blacklists operated by antivirus vendors and found that, for most prevalent malware families, the
union of all 15 public blacklists covered less than 20\% of the families' malicious domains. Li et
al.\ \cite{li2019teaLeaves} formalised a set of feed metrics --- volume, differential contribution,
exclusive contribution, latency, coverage and accuracy --- and applied them to a broad range of
public and commercial feeds; they found that while a few feeds overlap substantially, most do not,
that most indicators are unique to a single feed, and that these properties are stable over time.
Bouwman et al.\ \cite{bouwman2020differentcup} extended the comparison to two leading
\emph{commercial} threat intelligence vendors and found almost no overlap between them or with four
large open feeds: even for the 22 threat actors both vendors claim to track, indicator overlap
averaged only 2.5--4.0\%, and indicators that did overlap appeared in the second vendor's feed about
a month later on average.

On the aggregator side, Zhu et al.\ \cite{zhu2020labeldynamics} collected daily VirusTotal labels
for more than 14\,000 files from 65 engines over a year and showed that individual engine labels
change over time, that certain groups of engines are strongly correlated, and that threshold-based
aggregation of engine verdicts stabilises labels. Their setting is file scanning rather than IP
reputation, but the two observations carry over directly: the per-engine time series measured here
are expected to flip, and correlated engine groups (such as the BitDefender/G-Data pair in
\S\ref{sec:agreement}) must be accounted for before reading agreement as independent confirmation.

What distinguishes this thesis from that body of work is the reference point. Prior studies compare
feeds with each other, so they can measure \emph{relative} latency (one feed lists an indicator
before another) but not lag from the moment the indicator was known to be malicious. Anchoring every
tracked IP to the date of an independent TLScope classification allows lag to be measured from a
fixed origin, and following each IP for several weeks allows the same flags to be followed to
expiry.

\section{Malicious Domain and IP Reputation Systems}

\subsection{NOTOS}

Antonakakis et al.\ \cite{antonakakis2010notos} introduced NOTOS, a dynamic reputation system for
DNS that scores domains before they appear on static blacklists, using passive DNS data from ISP
resolvers and network and zone features (IP diversity, TTL volatility, registration age, registrar
identity). It achieves a 96.8\% true-positive rate at a 0.38\% false-positive rate and detects
malicious domains weeks to months before blacklists. The ``lead time before blacklists'' it reports
is the detector-side counterpart of the detection lag measured in this thesis.

\subsection{EXPOSURE}

Bilge et al.\ \cite{bilge2011exposure} extended passive DNS analysis to 15 time-, answer-, TTL- and
name-based features, achieving about 98\% accuracy on 100 billion DNS requests, and measured an
average malicious-domain lifespan of 1.2 days in live deployment. That figure concerns how long the
malicious \emph{infrastructure} stays active, whereas this thesis measures how long a \emph{provider}
keeps its flag raised; the two are complementary, and the gap between them is the interval in which
a feed reports on infrastructure that may already be gone.

\subsection{MANTIS}

Deniz et al.\ \cite{deniz2025mantis} presented MANTIS, a content-agnostic system for zero-day
malicious domain detection that builds a heterogeneous graph over domains, IPs, ASNs, nameservers
and registrars and applies a graph neural network to it. It reports 99.7\% precision and 86.9\%
recall and detects roughly 19\,000 new malicious domains per day, about five times as many as
VirusTotal alone. MANTIS, like TLScope, is evidence that infrastructure-level detectors find
substantially more than public feeds; this thesis measures how much of what one such detector finds
the feeds eventually catch up with, and when.

\section{Machine Learning for Network Security}

\subsection{Robustness Under Concept Drift}

Li et al.\ \cite{li2025conceptdrift} addressed concept drift in Windows malware detection using graph
neural networks on control-flow graphs and adversarial domain adaptation. Drift is relevant here
because the providers' own classifiers drift: the verdict an engine returns for the same IP can
change over time independently of the server's behaviour, which is one of the expiry mechanisms
discussed in \S\ref{sec:expiry}.

\subsection{Low-Quality Training Data}

Qing et al.\ \cite{qing2024rapier} presented RAPIER, a framework for encrypted malicious traffic
detection under up to 45\% label noise, built on the observation that malicious traffic is scattered
across feature space while benign traffic is concentrated. By analogy, an IP flagged by many
independent engines is a high-confidence indicator, while a single-engine flag is the noise-prone
case. In this dataset the single-engine case dominates (83.1\% of detected IPs,
\S\ref{sec:agreement}), so most provider-side confirmations of TLScope's verdicts rest on one
engine's judgement.

\subsection{Density-Based Robustness}

Tian et al.\ \cite{tian2025density} identified feature sparsity as a root cause of several failure
modes of AI-driven malware detectors, including poisoning, evasion and concept drift. The
TLScope-identified IPs that no provider ever flags are plausibly the sparse region of the providers'
own indicator space; this interpretation is offered as a hypothesis and is not tested here
(Chapter~\ref{ch:future_work}).

\section{Internet-Wide Scanning and Infrastructure Analysis}

\subsection{IPv6 Scanning with 6Sense}

Williams et al.\ \cite{williams2024sixsense} presented 6Sense, which uses reinforcement learning to
make IPv6 internet scanning tractable. This thesis tracks only IPv4; 6Sense establishes the
groundwork for an IPv6 extension, where public feed coverage may lag further still.

\subsection{BGP Routing Anomaly Detection (BEAM)}

Chen et al.\ \cite{chen2024beam} presented BEAM, a semantics-aware BGP anomaly detector built on
AS-level role embeddings. Its notion of role consistency suggests an analogous, untested hypothesis
here: that an IP's detection profile across providers is driven by the hosting role of its
autonomous system rather than by the engines themselves.

\subsection{Proxy-Based Abuse Detection (CalcuLatency)}

Ramesh et al.\ \cite{ramesh2024calcuLatency} presented CalcuLatency, which detects proxy-enabled
abuse from cross-layer latency measurements. Latency fingerprints could group TLScope-identified IPs
by hosting infrastructure, allowing a future study to test whether IPs on the same infrastructure
share detection lag --- the bulk listing events observed in \S\ref{sec:lag} suggest they do.

\section{DNS Infrastructure and Security}

Zhang et al.\ \cite{zhang2024staleglue} showed that 23.18\% of DNS glue records are outdated, and
Zhang et al.\ \cite{zhang2024resolverfuzz} found 23 vulnerabilities in mainstream DNS resolvers. Both
bear on the one provider in this study that is consulted through DNS: OpenPhish URLs are resolved to
IPs before matching, so stale or incorrect resolution can attribute a phishing indicator to the
wrong address (\S\ref{sec:openphish}).

% ============================================================
\chapter{Methodology}
\label{ch:methodology}
% ============================================================

\section{Overview}

NetScan follows a fixed population of TLScope-identified malicious IPs forward in time, querying
every provider for every IP at regular intervals and recording each provider's (and each VirusTotal
engine's) verdict on each sample. Detection lag, indicator lifetime and agreement are then derived
from those per-sample verdicts.

\section{Data Collection}

\subsection{Waves and Sub-Groups}
\label{sec:waves}

The tracked population is organised in \emph{waves}. Each wave consists of two or three
\emph{sub-groups}, and each sub-group is seeded by one dated TLScope output batch (a CSV with
columns \texttt{IP}, \texttt{Port} and \texttt{label}; only rows whose label is not
\texttt{benign} are tracked). The sub-groups of a wave start on consecutive days, so that the
collection tool queries one sub-group per day; each individual IP is therefore sampled once per
cycle of $\Delta$ days, not daily. Table~\ref{tab:waves} summarises the seven waves analysed.

\begin{table}[h]
  \centering
  \caption{Waves analysed in this thesis. $\Delta$ is the sampling interval of each IP; the window
    is the longest span from a sub-group's first to last sample. Wave 8's window is the analysed
    portion (\S\ref{sec:data_quality}).}
  \label{tab:waves}
  \small
  \setlength{\tabcolsep}{4.5pt}
  \begin{tabular}{lllrrrrc}
    \toprule
    \textbf{Wave} & \textbf{Sub-groups} & \textbf{First start} & $\boldsymbol{\Delta}$ \textbf{(d)}
      & \textbf{Samples/IP} & \textbf{Window (d)} & \textbf{IPs} & \textbf{Censys} \\
    \midrule
    2 & d, e, f & 2025-12-04 & 3 & 11.8 & 36 & 2\,062 & yes \\
    3 & g, h, i & 2026-01-13 & 3 &  9.7 & 33 & 2\,728 & yes \\
    4 & j, k, l & 2026-02-19 & 3 & 11.9 & 39 & 2\,474 & yes \\
    5 & m, n    & 2026-04-04 & 2 & 17.2 & 34 & 1\,596 & partly \\
    6 & o, p, q & 2026-05-12 & 3 & 11.8 & 36 & 2\,470 & no \\
    7 & r, s, t & 2026-06-23 & 3 & 12.3 & 36 & 2\,100 & no \\
    8 & u, v, w & 2026-08-04 & 3 & 10.6 & 39 & 2\,407 & no \\
    \midrule
    \textbf{Total} & & & & & & \textbf{15\,837} & \\
    \bottomrule
  \end{tabular}
\end{table}

\textbf{Alignment of tracking start with TLScope classification.} Tracking of each sub-group begins
on the date its TLScope batch was produced --- the batch directory name carries that date (e.g.\
\texttt{tlscope\_keramidas\_2026-08-05} is tracked from 2026-08-05) --- for every sub-group in
Table~\ref{tab:waves}. This alignment is what makes RQ1 answerable: ``days since the sub-group's
start'' is a direct proxy for ``days since TLScope identified the server''. Were the two decoupled,
the measured lag would describe only when this study began looking.

\subsection{Collection Schedule}
\label{sec:collection_schedule}

On each run \texttt{runfetch.py} queries every configured provider for every IP of the sub-group due
that day and writes a timestamped report CSV. Six of the seven waves sample each IP every three
days; wave 5 samples every two days, which the analysis takes into account through a per-sub-group
\texttt{config.json} override (Appendix~\ref{app:cli}). Lifetimes are computed in days using each
wave's own $\Delta$, so they are comparable across waves; detection lag is measured in days from the
sub-group start and inherits a resolution of $\Delta$.

\subsection{Provider Availability}
\label{sec:provider_availability}

All six providers were queried throughout, with one exception. Censys querying was removed from the
collection tool's active provider list (\texttt{runfetch.py}, commit \texttt{68d4048}, 2026-04-19)
and never re-enabled. Censys responses disappear from the wave-5 report files between 2026-04-15 and
2026-04-17, so Censys covers waves 2--4 (7\,264 IPs) fully, the first twelve days of wave 5, and
none of waves 6--8. Censys produced no detection in any wave-5 sample in which it was queried. No
rationale for the removal is recorded in the codebase; its effect on the results is addressed in
\S\ref{sec:threats_to_validity}.

\subsection{Provider Query Strategies}

Table~\ref{tab:provider_queries} summarises how each provider is queried and what counts as a
detection; Chapter~\ref{ch:providers} gives the details.

\begin{table}[h]
  \centering
  \caption{Provider query strategies and classification basis.}
  \label{tab:provider_queries}
  \begin{tabularx}{\linewidth}{lXX}
    \toprule
    \textbf{Provider} & \textbf{Query type} & \textbf{Detection} \\
    \midrule
    VirusTotal  & POST rescan, then GET analysis & Per engine: category \texttt{malicious} or \texttt{suspicious} \\
    AbuseIPDB   & GET check (30-day report window) & \texttt{abuseConfidenceScore} $>$ 15 \\
    Censys      & GET host record                 & A \texttt{c2} label on the host or a service \\
    ThreatFox   & POST IOC search (exact match)   & \texttt{query\_status} $\neq$ \texttt{no\_result} \\
    Cinsscore   & Local lookup in downloaded list & IP present in CINS Army list \\
    OpenPhish   & DNS resolve + list lookup       & IP present among resolved feed hosts \\
    \bottomrule
  \end{tabularx}
\end{table}

\subsection{Rate Limiting}
\label{sec:rate_limiting}

VirusTotal enforces hourly and daily API quotas. A background \texttt{quota\_worker} task monitors the
remaining quota and cancels all fetch tasks when the daily limit is reached, and a wait of 32 seconds
(\texttt{WAITING\_TIME}) is applied between successive IPs. Cinsscore and OpenPhish are list-based:
the full list is downloaded once at startup and all IPs are checked locally.

\section{Measurement Definitions}
\label{sec:notation}

For a tracked IP $i$ in a sub-group with start date $d_0$, let $d_1 < d_2 < \dots < d_n$ be the
dates on which $i$ was sampled ($d_1 = d_0$ in the normal case, consecutive dates $\Delta$ apart),
and for each \emph{engine} $e$ --- a VirusTotal engine or one of the five other providers --- let
$x_t(i,e) \in \{0,1\}$ indicate whether $e$ classified $i$ as malicious or suspicious on sample $t$.
IP $i$ is \emph{detected} on sample $t$ if $x_t(i,e) = 1$ for at least one engine.

\textbf{Detection lag.} If $i$ is detected on its first sample and $d_1 = d_0$, it has \emph{zero
lag}: it was already flagged when TLScope identified it (to a resolution of $\Delta$). If it is
not detected on the first sample but is detected later, its lag is $d_\tau - d_0$ days, where
$\tau$ is the first detected sample. If $i$'s first recorded sample falls after $d_0$ --- because the
IP's first query failed or its batch file entered the schedule late --- and it is detected, its lag
is \emph{undetermined}. If $i$ is never detected, its lag is right-censored at the window length.

\textbf{Indicator lifetime.} An indicator is considered expired after two consecutive non-detections
(\texttt{EXPIRED\_AFTER} $= 2$), so a single isolated miss does not end it. Each sample from the
first detection $\tau(i,e)$ onwards contributes $\Delta$ days to the lifetime until the second
consecutive miss at sample $t^*$:
\[
  L(i,e) = \bigl(t^*(i,e) - \tau(i,e)\bigr)\cdot\Delta ,
\]
i.e.\ the detected samples, any single intervening misses, and the first miss of the expiring pair.
If the indicator never expires within the window, every sample from $\tau$ to $n$ counts,
$L = (n - \tau + 1)\cdot\Delta$, and the lifetime is \emph{right-censored}: the true value is at
least $L$. Only the first detection episode is measured; a re-detection after expiry is not counted
as a new indicator. The \emph{censored share} of an engine is the fraction of its indicators still
active on the last sample.

\textbf{Peak detection share.} For each engine and each sample index, the share of tracked IPs it
flags among those for which it returned a verdict; an engine's \emph{peak} is the maximum of this
series over the window. The five engines with the highest peak in each wave are reported in detail
(\texttt{-t 5}).

\textbf{Agreement.} Let $D_e$ be the set of IPs engine $e$ ever flags in a wave. An IP's agreement
count is the number of engines whose $D_e$ contains it; pairwise overlap is the Jaccard index
$J(a,b) = |D_a \cap D_b| \,/\, |D_a \cup D_b|$.

\textbf{Independence baseline.} To ask whether engines overlap more or less than their individual
coverage alone implies, each engine is modelled as flagging each of the wave's $N_{\text{wave}}$
tracked IPs independently with probability $p_e = |D_e| / N_{\text{wave}}$. The number of engines
flagging a given IP then follows a Poisson binomial distribution, from which follow the expected
share of IPs flagged by at least one engine and, among those, the expected share flagged by two or
more. Both are compared with the observed values (\S\ref{sec:agreement_vs_chance}).

\section{Data Quality and Corrections}
\label{sec:data_quality}

Because every result rests on the collected report files, the collection itself was audited before
analysis. Four issues were found.

\textbf{Pilot wave excluded.} The first wave (October--November 2025, 2\,591 IPs) was collected while
the collection tool was still being developed. Its report files do not share one schema --- some
lack the ThreatFox, Cinsscore and OpenPhish columns because those providers were added to the tool during the wave
--- and one run was made with a defective ThreatFox classifier: the version of
\texttt{process\_tf\_report} committed on 2025-10-29 marked an IP as malicious whenever ThreatFox's
response contained a non-empty \texttt{data} field, which includes ThreatFox's own ``no result''
message. All 139 ThreatFox ``detections'' in that wave come from this one run and every one of them
carries a \texttt{no\_result} response. The classifier was fixed on 2025-10-31, before any later
wave began, and no later wave contains a ThreatFox detection with a \texttt{no\_result} response.
Because provider coverage within the pilot wave is not uniform, it is excluded from all results
rather than patched.

\textbf{Censys deactivation.} As described in \S\ref{sec:provider_availability}, Censys covers only
waves 2--4 and the start of wave 5. This is detectable in the data itself: Censys fields vanish from
the per-wave field inventory (\texttt{group\_available\_fields.json}) from wave 6 onwards.

\textbf{Incomplete VirusTotal responses.} A sample row whose VirusTotal analysis could not be
retrieved carries fewer than 50 engine verdicts. Between 0.5\% and 3.0\% of sample rows per wave are
in this state (3.0\% in wave 8, at most 1.8\% elsewhere); they are kept, since the remaining providers' verdicts in the same row are valid, but
they can only under-count detections.

\textbf{Wave 8 cut-off.} Collection for wave 8 continued after its analysis was frozen. To keep wave
8 comparable with the other waves, it is analysed up to its sub-groups' samples of 2026-09-10 to
2026-09-12 (a 39-day window); the effect of the longer collection is reported separately as a
sensitivity check (\S\ref{sec:window_sensitivity}).

All figures in Chapters~\ref{ch:results} and~\ref{ch:discussion} were computed from the raw report
files of waves 2--8 by an analysis that reproduces \texttt{runanalysis.py}'s per-wave output exactly
(total, never-detected, lag and agreement counts and per-engine lifetimes) and additionally records
the per-IP distributions that the standard output summarises.

\section{Ethical Considerations}

Every provider is queried through its own documented public API or download endpoint, within the rate
limits each enforces (\S\ref{sec:rate_limiting}); no access controls are bypassed, and no
authentication beyond a standard per-provider API key (Appendix~\ref{app:config}) is used. NetScan
does not itself connect to the tracked malicious IPs: five of the six providers return
already-collected reputation data without contacting the target. The exception is VirusTotal, whose
rescan endpoint causes VirusTotal's own infrastructure --- not NetScan --- to probe the target. No
traffic to or from the tracked infrastructure is captured or decrypted. The tracked indicators are
server-side infrastructure addresses from an already published upstream classification, not
end-user or personally identifiable data; where individual IPs are shown in this thesis, their last
two octets are masked.

% ============================================================
\chapter{Threat Intelligence Providers}
\label{ch:providers}
% ============================================================

\section{VirusTotal}

VirusTotal aggregates the verdicts of 90--100 third-party blocklist, reputation and security
engines for an IP address. For each IP, NetScan first triggers a fresh analysis via
\texttt{POST /api/v3/ip\_addresses/\{IP\}/analyse} and then polls it until
\texttt{status == completed}. Each engine's verdict is recorded individually, and an engine counts as
detecting the IP when its category is \texttt{malicious} or \texttt{suspicious}. VirusTotal is
therefore analysed as a collection of engines rather than as a single provider. The worker backs off
on \texttt{Too Many Requests} (64 s) and \texttt{QuotaExceededError} (128 s).

\section{AbuseIPDB}

AbuseIPDB is a community-driven database of reported IPs. A single GET request with a 30-day window
(\texttt{maxAgeInDays = 30}) returns the \texttt{abuseConfidenceScore} (0--100), to which NetScan
applies thresholds: $>$25 malicious, $>$15 suspicious, otherwise harmless. Because the score only
reflects reports from the last 30 days, an AbuseIPDB flag decays on its own once reports stop
arriving, which bounds its lifetime by construction.

\section{Censys}

Censys continuously scans the public IPv4 address space and maintains per-host labels including
\texttt{c2} for known command-and-control infrastructure. NetScan checks for \texttt{c2} labels at
both the host and the individual service level; either counts as malicious. Censys was queried only
in waves 2--4 and the beginning of wave 5 (\S\ref{sec:provider_availability}).

\section{ThreatFox}

ThreatFox (abuse.ch) is a free community platform for sharing malware IOCs. A POST request with
\texttt{exact\_match: true} is issued for the IP; any \texttt{query\_status} other than
\texttt{no\_result} is classified as malicious.

\section{Cinsscore}

The CINS Army list is a public text file of IPs with poor reputation, regenerated regularly. NetScan
downloads it once per run and checks all IPs locally.

\section{OpenPhish}
\label{sec:openphish}

OpenPhish publishes a feed of phishing URLs. Because the feed contains URLs rather than IPs, NetScan
extracts the hostnames and resolves them to IPs via DNS at the start of each run; an IP found among
the resolved addresses is classified as malicious. Resolution introduces noise: domains that rotate
IPs (fast-flux, CDN-backed phishing) produce transient associations, and stale DNS glue records
\cite{zhang2024staleglue} are an additional source of incorrect attribution.

% ============================================================
\chapter{Implementation}
\label{ch:implementation}
% ============================================================

\section{Repository Structure}

\begin{table}[h]
  \centering
  \caption{NetScan repository files and their purpose.}
  \label{tab:repo}
  \begin{tabularx}{\linewidth}{lX}
    \toprule
    \textbf{File} & \textbf{Purpose} \\
    \midrule
    \texttt{runfetch.py}       & Asynchronous data collection from all providers \\
    \texttt{runanalysis.py}    & Analysis pipeline and graph generation \\
    \texttt{vendors.py}        & Provider API client functions and report processors \\
    \texttt{available\_fields.py} & Utility to inspect the CSV output field schema \\
    \texttt{.env\_ex}          & Example environment / API key configuration \\
    \bottomrule
  \end{tabularx}
\end{table}

\section{Fetch Pipeline}

VirusTotal's per-IP analysis is asynchronous on VirusTotal's side: a submitted rescan must be polled
until its status becomes \texttt{completed}, and each poll waits a multiple of
\texttt{WAITING\_TIME}. An earlier single-pool design queried every provider for a given IP from the
same worker, which then sat blocked on VirusTotal's polling loop even though the other providers had
already answered. The pipeline is therefore split into two cooperating pools connected by an
internal queue:

\begin{enumerate}
  \item At startup, the Cinsscore and OpenPhish lists are downloaded once and held in memory.
  \item For each input CSV in the folder, a new output file is created.
  \item A pool of \texttt{-w} \textbf{producer} workers pulls IPs from the input queue. For each IP,
    a producer submits the VirusTotal rescan \emph{without waiting for it to complete}, then queries
    every other configured provider. The partial result --- including the VirusTotal analysis id ---
    is handed to a second, VirusTotal-result queue.
  \item A pool of \texttt{-vw} \textbf{consumer} workers polls each VirusTotal analysis until it
    completes, merges the result into the row, and writes the row to the output CSV.
  \item A separate \texttt{quota\_worker} task monitors VirusTotal's daily quota and cancels every
    producer and consumer if it is exhausted.
  \item Producers sleep 32 seconds between IPs to respect per-provider rate limits; the VirusTotal
    polling loop in each consumer applies its own backoff.
\end{enumerate}

This decouples the dominant source of per-IP latency from the retrieval of every other provider's
data. File I/O is performed asynchronously via \texttt{anyio} so that large reads and writes do not
block the event loop that drives the provider queries.

\section{Fetch Output Format}

Each report CSV contains one row per tracked \texttt{IP}:\texttt{Port} with the columns
\texttt{total\_votes} (aggregate vote counts), \texttt{report} (per-engine classification, VirusTotal
engines and the other providers alike) and the raw provider responses (\texttt{vt\_response},
\texttt{ai\_response}, \texttt{tf\_response}, \texttt{cb\_response}, \texttt{op\_response}, and
\texttt{cs\_response} while Censys was active). Keeping the raw responses is what made the
data-quality audit in \S\ref{sec:data_quality} possible.

\section{Analysis Output}

Running \texttt{runanalysis.py} on a wave folder writes the following files to an output
subdirectory (\texttt{group\_analysis\_out\_dir/} by default) inside that folder:

\begin{table}[h]
  \centering
  \caption{Files produced by the analysis pipeline.}
  \label{tab:analysis_output}
  \begin{tabularx}{\linewidth}{lX}
    \toprule
    \textbf{File} & \textbf{Content} \\
    \midrule
    \texttt{out\_logs.txt}       & Run arguments and all computed statistics (Total Stats, Report Stats) \\
    \texttt{all\_zero.csv}       & IPs never detected by any provider \\
    \texttt{non\_zero.csv}       & IPs detected by at least one provider at least once \\
    \texttt{ip\_port\_series.csv}  & Per-host vote history (\texttt{IP}, \texttt{Port}, \texttt{Group}, \texttt{Date}, votes) \\
    \texttt{total\_series.csv}   & Aggregate/diff/percentage series, one row per collection cycle \\
    \texttt{engine\_series.csv}  & Per-engine count/percentage series, one row per engine per cycle \\
    \texttt{engine\_summary.csv} & Per-engine average lifetime, peak detection share and
      ever-detected IP count \\
    \texttt{group\_available\_fields.json} & Merged schema of every field across all provider
      responses in the wave's report files, for schema inspection without hand-parsing raw CSVs \\
    \texttt{figure\_*.png}       & Generated graphs (aggregate, diff, top-engine detail) \\
    \bottomrule
  \end{tabularx}
\end{table}

% ============================================================
\chapter{Results}
\label{ch:results}
% ============================================================

This chapter reports results for waves 2--8 (15\,837 IPs, December 2025 -- September 2026). The
pilot wave is excluded throughout (\S\ref{sec:data_quality}). Methodological caveats are collected in
\S\ref{sec:threats_to_validity} rather than repeated at each table.

\section{Worked Examples}
\label{sec:examples}

Before the aggregates, Table~\ref{tab:examples} shows five real IPs from wave 4, one per typical
pattern, to make the measurement definitions of \S\ref{sec:notation} concrete. Each cell is the
number of engines flagging the IP on that sample.

\begin{table}[h]
  \centering
  \caption{Five wave-4 IPs sampled every 3 days (last two octets masked). Cells give the number of
    engines flagging the IP on each sample; $\cdot$ is no detection. Lag and lifetime follow
    \S\ref{sec:notation}; ``$\geq$'' marks a lifetime still running when observation ended.}
  \label{tab:examples}
  \footnotesize
  \setlength{\tabcolsep}{2.6pt}
  \begin{tabular}{l*{14}{c}l}
    \toprule
    & \multicolumn{14}{c}{\textbf{Day since sub-group start}} & \\
    \cmidrule(lr){2-15}
    \textbf{IP} & 0 & 3 & 6 & 9 & 12 & 15 & 18 & 21 & 24 & 27 & 30 & 33 & 36 & 39 & \textbf{Reading} \\
    \midrule
    (a) 206.237.x.x & 11 & 12 & 12 & 11 & 11 & 11 & 13 & 14 & 14 & 14 & 14 & 14 & 14 & 14 & lag 0; 15 engines \\
    (b) 190.5.x.x   & 1 & 1 & 1 & 1 & 1 & 1 & 1 & 1 & $\cdot$ & $\cdot$ & $\cdot$ & $\cdot$ & $\cdot$ & $\cdot$ & lag 0; $L = 27$ d \\
    (c) 193.135.x.x & 1 & 1 & $\cdot$ & 1 & 1 & 1 & 1 & 1 & 1 & 1 & 1 & 1 & 1 & 1 & lag 0; $L \geq 42$ d \\
    (d) 43.174.x.x  & $\cdot$ & $\cdot$ & $\cdot$ & $\cdot$ & $\cdot$ & $\cdot$ & $\cdot$ & $\cdot$ & $\cdot$ & $\cdot$ & 1 & 1 & 1 & 1 & lag 30 d; $L \geq 12$ d \\
    (e) 43.159.x.x  & $\cdot$ & $\cdot$ & $\cdot$ & $\cdot$ & $\cdot$ & $\cdot$ & $\cdot$ & $\cdot$ & $\cdot$ & $\cdot$ & $\cdot$ & $\cdot$ & $\cdot$ & $\cdot$ & never detected \\
    \bottomrule
  \end{tabular}
\end{table}

IP (a) is the rare consensus case: already flagged by eleven engines on the day TLScope identified
it, it accumulates fifteen distinct engines over the window --- among them Censys, ThreatFox and both
BitDefender and G-Data --- and none of them lets go. IP (b) is flagged by Criminal IP alone, which
holds the flag for eight samples and then drops it; the second consecutive miss on day 27 confirms
the expiry, giving $L = 9 \times 3 = 27$ days. IP (c), also Criminal IP alone, shows why a single
miss is not treated as expiry: the flag disappears on day 6 and returns on day 9, and the indicator
is still active at the end of the window. IP (d) is a late detection: alphaMountain.ai flags it 30
days after TLScope did. IP (e), the most common case by far, is never flagged by anyone; 1\,111 of
wave 4's 2\,474 IPs had a complete sample series and no detection at all.

\section{Coverage}
\label{sec:coverage}

Table~\ref{tab:detection_overview} and Figure~\ref{fig:outcome} give, for each wave, how many
TLScope-identified IPs were ever flagged by at least one engine during the window.

\begin{table}[h]
  \centering
  \caption{Detection outcome per wave. ``Lag'' columns split the ever-detected IPs by when they were
    first flagged (\S\ref{sec:notation}).}
  \label{tab:detection_overview}
  \small
  \begin{tabular}{lrrrrrrr}
    \toprule
    & & \multicolumn{2}{c}{\textbf{Never detected}} & & \multicolumn{3}{c}{\textbf{Ever detected, by lag}} \\
    \cmidrule(lr){3-4}\cmidrule(lr){6-8}
    \textbf{Wave} & \textbf{IPs} & \textbf{IPs} & \textbf{\%} & \textbf{Ever} & \textbf{Zero} & \textbf{Later} & \textbf{Undet.} \\
    \midrule
    2 & 2\,062 & 1\,507 & 73.1 & 555 & 215 & 339 &   1 \\
    3 & 2\,728 & 2\,188 & 80.2 & 540 & 299 & 109 & 132 \\
    4 & 2\,474 & 2\,144 & 86.7 & 330 & 228 &  92 &  10 \\
    5 & 1\,596 & 1\,390 & 87.1 & 206 & 145 &  56 &   5 \\
    6 & 2\,470 & 2\,079 & 84.2 & 391 & 177 & 176 &  38 \\
    7 & 2\,100 & 1\,772 & 84.4 & 328 & 141 & 159 &  28 \\
    8 & 2\,407 & 1\,987 & 82.6 & 420 & 116 & 209 &  95 \\
    \midrule
    \textbf{Total} & \textbf{15\,837} & \textbf{13\,067} & \textbf{82.5} & \textbf{2\,770}
      & \textbf{1\,321} & \textbf{1\,140} & \textbf{309} \\
    \bottomrule
  \end{tabular}
\end{table}

\begin{figure}[h]
  \centering
  \includegraphics[width=0.9\linewidth]{outcome_by_wave.png}
  \caption{Detection outcome per wave as a share of the wave's tracked IPs.}
  \label{fig:outcome}
\end{figure}

Across the seven waves, \textbf{82.5\% of TLScope-identified malicious IPs (13\,067 of 15\,837)
were never flagged by any provider} during their 33--39-day window. This is the headline coverage
result: the large majority of servers that active TLS fingerprinting identifies as malicious do not
surface in public threat intelligence within five weeks, if at all.

The rate ranges from 73.1\% to 87.1\% across waves. It is higher in the later waves, but with seven
waves the correlation with wave order is not significant (Pearson $r = 0.55$, $p = 0.20$), and
neither is the difference between the waves in which Censys was queried (2--4, mean 80.0\%) and
those in which it was not (5--8, mean 84.5\%; Welch $t = -1.32$, $p = 0.24$). The data therefore
show no reliable trend over the study period --- and, in particular, no sign that coverage is
improving.

\section{Detection Lag}
\label{sec:lag}

Of the 2\,770 detected IPs, \textbf{1\,321 (47.7\%) were already flagged on the first sample}, i.e.\
by the time TLScope identified them, and \textbf{1\,140 (41.2\%) were first flagged on a later
sample}. The remaining 309 (11.2\%) have an undetermined lag because their first recorded sample
fell after their sub-group's start date (\S\ref{sec:notation}); most of them are in waves 3 and 8.

The share of detected IPs that were already flagged varies widely, from 27.6\% (wave 8) to 70.4\%
(wave 5), with no trend across waves (Pearson $r = -0.40$, $p = 0.38$).

For the 1\,140 late detections, Figure~\ref{fig:lag} shows the lag in days. The median is
\textbf{15 days} (interquartile range 6--18); 37.4\% are flagged within 9 days and 71.1\% within 15,
but a long tail reaches the end of the window. The distribution is not smooth, and the reason is
visible in the data: late detections often arrive in bulk. In wave 2, 262 of the 339 late
detections happened on the same sample --- day 15 --- when Criminal IP added 266 wave-2 IPs to its
listings at once and then kept them (Figure~\ref{fig:share_all}). Excluding wave 2, the median lag of
late detections is 12 days (interquartile range 6--24).

\begin{figure}[h]
  \centering
  \includegraphics[width=0.85\linewidth]{lag_days_hist.png}
  \caption{Days from sub-group start to first detection, for the 1\,140 IPs first flagged after
    their first sample (waves 2--8 pooled, 3-day bins).}
  \label{fig:lag}
\end{figure}

\section{Indicator Lifetime}
\label{sec:lifetime}

Lifetimes are computed for every engine that ever flags a tracked IP: each VirusTotal engine
separately (e.g.\ \emph{Criminal IP}, \emph{SOCRadar}, \emph{BitDefender}), plus the five other
providers. Before any lifetime is read, one property of the data dominates: \textbf{most flags never
expire within the window}. Of the 3\,945 engine--IP indicators across waves 2--8, 2\,550 (64.6\%)
are still active on the last sample. For those, the measured lifetime is the time from first
detection to the end of observation, and the true lifetime is longer.

Two engines are among the five with the highest peak detection share in every wave:
\emph{Criminal IP} and \emph{alphaMountain.ai}. Table~\ref{tab:top_engine} tracks both.

\begin{table}[h]
  \centering
  \caption{Average indicator lifetime, censored share (indicators still active at the end of the
    window) and peak detection share for the two engines among the five highest-peak engines in
    every wave.}
  \label{tab:top_engine}
  \small
  \begin{tabular}{lrrrrrr}
    \toprule
    & \multicolumn{3}{c}{\textbf{Criminal IP}} & \multicolumn{3}{c}{\textbf{alphaMountain.ai}} \\
    \cmidrule(lr){2-4}\cmidrule(lr){5-7}
    \textbf{Wave} & \textbf{Lifetime (d)} & \textbf{Censored} & \textbf{Peak}
                  & \textbf{Lifetime (d)} & \textbf{Censored} & \textbf{Peak} \\
    \midrule
    2 & 22.1 & 86\% & 18.7\% & 33.6 & 100\% & 1.7\% \\
    3 & 25.5 & 78\% & 12.8\% & 24.6 &  96\% & 1.7\% \\
    4 & 32.0 & 80\% &  8.0\% & 20.0 & 100\% & 3.0\% \\
    \midrule
    5 & 11.1 & 15\% &  5.0\% & 37.0 &  98\% & 3.0\% \\
    6 & 11.3 &  6\% &  5.5\% & 33.3 & 100\% & 1.5\% \\
    7 &  8.4 &  8\% &  4.0\% & 36.0 &  98\% & 2.3\% \\
    8 &  8.1 & 29\% &  4.3\% & 30.4 & 100\% & 2.3\% \\
    \bottomrule
  \end{tabular}
\end{table}

The two engines show different kinds of persistence. \emph{alphaMountain.ai} essentially never
withdraws a flag: in every wave at least 96\% of its indicators are still active at the end, so its
20--37-day average lifetimes are set by how far into the window it first flagged each IP, not by any
expiry. Its coverage is narrow throughout (peak share at most 3.0\%).

\emph{Criminal IP} is the broadest single engine in every wave, and its behaviour changes
part-way through the study. In waves 2--4 (December 2025 -- March 2026) it behaved like
alphaMountain.ai: 78--86\% of its flags were still in place at the end of the window and average
lifetimes were 22--32 days. From wave 5 (April 2026) onwards, only 6--29\% of its flags survive to
the end and average lifetimes fall to 8--11 days --- it now withdraws most listings within one to two
weeks. Its peak share fell at the same time, from 8.0--18.7\% to 4.0--5.5\%. The drop in wave 5's
detection share around day 10 (Figure~\ref{fig:share_all}) is this behaviour at work: Criminal IP
released 35 of its 61 flagged wave-5 IPs between two consecutive samples. Whether the change reflects
a retention-policy change at Criminal IP or a change in the infrastructure TLScope surfaced cannot be
told apart from this data; the timing, consistent across four consecutive waves, points to the
former.

Beyond these two, the top five turns over from wave to wave. \emph{SOCRadar} appears in six of the
seven waves (2--7; average lifetime 24.1--32.4 days); \emph{AbuseIPDB} (waves 2, 6, 7) and
\emph{Gridinsoft} (6--8) in three each; \emph{Fortinet} (3, 4) and \emph{ESET} (5, 8) in two each;
and \emph{CRDF}, \emph{ThreatFox}, \emph{MalwareURL}, \emph{Seclookup} and \emph{Chong Lua Dao} in one
wave apiece. Twelve distinct engines fill the 35 top-five slots, so which engines contribute
detections is itself unstable over a nine-month study.

\section{Provider-Level Coverage}
\label{sec:provider_coverage}

The analysis above works at the level of VirusTotal's individual engines, because those are what
reach the highest detection shares. Table~\ref{tab:provider_coverage} asks the provider-level
question for the five providers queried outside VirusTotal.

\begin{table}[h]
  \centering
  \caption{IPs ever flagged by each standalone provider, waves 2--8. Censys's percentage is of the
    7\,264 IPs in waves 2--4, where it was queried throughout (\S\ref{sec:provider_availability}); the
    others are of all 15\,837 IPs. Lifetimes are averages over each provider's indicators and are
    right-censored like those in Table~\ref{tab:top_engine}.}
  \label{tab:provider_coverage}
  \small
  \begin{tabular}{lrrr}
    \toprule
    \textbf{Provider} & \textbf{Ever-flagged IPs} & \textbf{\% of population} & \textbf{Avg.\ lifetime (d)} \\
    \midrule
    AbuseIPDB  & 170 & 1.07\%                  & 24.8 \\
    ThreatFox  & 120 & 0.76\%                  & 7.0  \\
    Censys     & 16  & 0.22\% (waves 2--4)     & 3.4  \\
    OpenPhish  & 2   & 0.01\%                  & 7.5  \\
    Cinsscore  & 0   & 0.00\%                  & ---  \\
    \bottomrule
  \end{tabular}
\end{table}

\textbf{Cinsscore never flagged any of the 15\,837 tracked IPs}, and OpenPhish flagged two, both in
wave 3. OpenPhish's near-zero coverage is best read as a construct mismatch rather than a failure:
it is a phishing-\emph{URL} feed, while TLScope's ground truth targets TLS-fingerprinted C2 and
botnet servers. ThreatFox and AbuseIPDB are genuine but modest contributors. ThreatFox's detections
are concentrated (99 of its 120 are in wave 3) and short-lived: most come from a single sample on
day 24 of wave 3, when it briefly listed 77 wave-3 IPs and withdrew 70 of them three days later.
AbuseIPDB's longer lifetimes are consistent with its 30-day reporting window. Censys's 16 detections
come from waves 3 and 4 only.

Taken together, the five standalone providers flag 291 of the 2\,770 detected IPs (10.5\%), and only
125 (4.5\%) are flagged by one of them and by no VirusTotal engine. \textbf{VirusTotal's engines
alone account for 2\,645 of the 2\,770 detected IPs (95.5\%)}, so in this dataset the aggregator is
the source of almost all coverage, and the standalone providers add little that its engines do not
already report.

\section{Inter-Provider Agreement}
\label{sec:agreement}

Table~\ref{tab:agreement} counts, for every detected IP, how many engines ever flagged it.

\begin{table}[h]
  \centering
  \caption{Number of engines ever flagging each detected IP, waves 2--8.}
  \label{tab:agreement}
  \begin{tabular}{lrr}
    \toprule
    \textbf{Engines} & \textbf{IPs} & \textbf{\% of detected} \\
    \midrule
    Exactly 1  & 2\,303 & 83.1\% \\
    Exactly 2  & 271    & 9.8\%  \\
    Exactly 3  & 86     & 3.1\%  \\
    4 or more  & 110    & 4.0\%  \\
    \midrule
    \textbf{Total} & \textbf{2\,770} & \textbf{100.0\%} \\
    \bottomrule
  \end{tabular}
\end{table}

Agreement is low in absolute terms: 83.1\% of detected IPs are flagged by exactly one engine, and
only 4.0\% by four or more (up to 18 for the most-flagged IP). One pair is a clear exception: the
Jaccard index between \emph{BitDefender} and \emph{G-Data} is exactly 1.0 in all seven waves --- every
IP one flags, the other flags too. This is consistent with G-Data's products using Bitdefender's
scanning engine, so the pair should be read as one source reported twice rather than as two
independent confirmations, as Zhu et al.\ found for correlated engine groups in file scanning
\cite{zhu2020labeldynamics}.

\subsection{Agreement Versus an Independence Baseline}
\label{sec:agreement_vs_chance}

Low overlap has two candidate explanations: engines may each flag so few IPs that they rarely
coincide even by chance (\emph{sparsity}), or they may systematically flag different IPs from one
another (\emph{specialisation}). The independence baseline of \S\ref{sec:notation} separates the
two: if engines specialise, the observed multi-engine share should fall \emph{below} what
independent engines with the same individual rates would produce; if engines tend to agree on a
common core of IPs, it should fall above. Table~\ref{tab:agreement_chance} gives both comparisons
per wave.

\begin{table}[h]
  \centering
  \caption{Observed values against an independence model in which each engine flags each tracked IP
    independently at its own observed rate (\S\ref{sec:notation}). ``Ever flagged'' is the share of
    the wave's IPs flagged by at least one engine; ``2+ engines'' is the share of flagged IPs
    flagged by two or more. ``Engines'' counts engines with at least one detection in the wave.}
  \label{tab:agreement_chance}
  \small
  \begin{tabular}{lrrrrrr}
    \toprule
    & & \multicolumn{2}{c}{\textbf{Ever flagged}} & \multicolumn{3}{c}{\textbf{2+ engines}} \\
    \cmidrule(lr){3-4}\cmidrule(lr){5-7}
    \textbf{Wave} & \textbf{Engines} & \textbf{Obs.} & \textbf{Indep.} & \textbf{Obs.} & \textbf{Indep.} & \textbf{Ratio} \\
    \midrule
    2 & 26 & 26.9\% & 34.1\% & 25.2\% & 13.9\% & 1.81 \\
    3 & 36 & 19.8\% & 29.1\% & 18.9\% & 12.9\% & 1.46 \\
    4 & 31 & 13.3\% & 17.5\% & 13.3\% &  6.9\% & 1.92 \\
    5 & 23 & 12.9\% & 14.1\% &  8.7\% &  4.9\% & 1.79 \\
    6 & 33 & 15.8\% & 18.8\% & 11.3\% &  6.6\% & 1.71 \\
    7 & 33 & 15.6\% & 21.2\% & 19.2\% &  9.2\% & 2.10 \\
    8 & 31 & 17.4\% & 21.1\% & 13.3\% &  7.9\% & 1.69 \\
    \bottomrule
  \end{tabular}
\end{table}

In every wave the observed multi-engine share is \emph{higher} than independence predicts, by a
factor of 1.46--2.10, and correspondingly fewer IPs are flagged at all than independent engines would
reach. Engines therefore overlap \emph{more} than chance, not less: there is a core of IPs that
several engines recognise, and the large single-engine share in Table~\ref{tab:agreement} arises
because each engine individually flags only a small fraction of the population. The low agreement is
a consequence of sparsity, not evidence of specialisation. The independence model is a
simplification --- it ignores, for example, that IPs on shared hosting may be flagged together --- but
such effects push in the same direction as the observed excess and do not change its sign.

\section{Temporal Evolution}
\label{sec:temporal}

Figure~\ref{fig:share_all} shows, for every wave, the share of tracked IPs flagged by at least one
engine on each sample. In four of the seven waves (4 and 6--8) the share stays within a band of a few
percentage points for the whole window, with no sustained climb: coverage reaches a rough
equilibrium early rather than catching up over time. The two departures are the events identified
above --- Criminal IP's bulk listing on day 15 of wave 2, which more than doubles that wave's share
and holds, and its bulk release around day 10 of wave 5 --- plus ThreatFox's one-sample burst on
day 24 of wave 3.

\begin{figure}[h]
  \centering
  \includegraphics[width=\linewidth]{detection_share_all_waves.png}
  \caption{Share of each wave's tracked IPs flagged by at least one engine, per sample.}
  \label{fig:share_all}
\end{figure}

The equilibrium is not static. Figure~\ref{fig:wave4_diff} decomposes wave 4's aggregate vote count
(the number of engine flags across all its IPs) into gains and losses between consecutive samples.
The flagged share stays between 8.9\% and 12.4\% for all 14 samples, yet every step contains both new
flags and withdrawn ones: the net change ranges from $-14$ to $+22$ votes per cycle, while gains of
8--33 and losses of 7--28 votes occur at every step. Coverage is stable in level but continuously
churning in composition.

\begin{figure}[h]
  \centering
  \includegraphics[width=0.85\linewidth]{wave4_diff.png}
  \caption{Net, positive and negative components of the step-to-step change in aggregate engine
    flags for wave 4.}
  \label{fig:wave4_diff}
\end{figure}

\section{Sensitivity to Window Length}
\label{sec:window_sensitivity}

Collection for wave 8 continued beyond the window analysed above, to about 60 days. Analysing the
same 2\,407 IPs over the full collection lowers the never-detected rate from 82.6\% (39 days) to
68.2\% (60 days): 1\,641 instead of 1\,987 IPs remain unflagged, and the number of late detections
rises from 209 to 552. Coverage therefore keeps growing after five weeks, and the 82.5\% headline
figure is specific to a 33--39-day window. It is a measure of how much public feeds cover
\emph{within the first five weeks} after an independent detector has found a server, not of how
much they eventually cover.

% ============================================================
\chapter{Discussion}
\label{ch:discussion}
% ============================================================

\section{Answers to the Research Questions}
\label{sec:rq_answers}

A precondition applies to all three answers: lag, lifetime and agreement can only be measured on the
17.5\% of tracked IPs that some engine ever flagged. The remaining 82.5\% produce no detection event
within the window, and that censoring is itself the single largest finding.

\textbf{RQ1 (detection lag) --- bimodal: already known, or one to three weeks later.} Of the
detected IPs, 47.7\% were flagged on the first sample and 41.2\% later, with the rest undetermined
(\S\ref{sec:lag}). Late detections took a median of 15 days (interquartile range 6--18) and often
arrived in bulk, as when Criminal IP listed 266 wave-2 IPs on one day. The already-flagged share
swings from 27.6\% to 70.4\% between waves with no trend, which points to the composition of each
batch rather than to a fixed ecosystem latency. Lag is resolved only to the sampling interval
(three days, two in wave 5).

\textbf{RQ2 (indicator persistence) --- mostly longer than the window, with one engine changing
policy.} Nearly two thirds of all detections (64.6\%) are still in place when observation ends, so
measured lifetimes are lower bounds. Among the engines present in every wave, alphaMountain.ai
withdraws almost nothing within the window, while Criminal IP --- the broadest engine --- switched in
April 2026 from holding flags for the whole window (22--32-day average lifetimes, 78--86\% still
active) to withdrawing most within one to two weeks (8--11 days, 6--29\% still active). Among
standalone providers, AbuseIPDB flags last longest (24.8 days on average, bounded by its 30-day
reporting window) and ThreatFox and Censys flags shortest (7.0 and 3.4 days). Persistence is
therefore a property of each engine's retention practice, and that practice can change without
notice.

\textbf{RQ3 (inter-provider agreement) --- low, because each engine sees little, not because they
see different things.} 83.1\% of detected IPs are flagged by exactly one engine and 4.0\% by four or
more. Compared with an independence model, engines co-detect 1.46--2.10 times more often than
chance in every wave, so the low overlap is explained by sparse individual coverage, not by
specialisation. BitDefender and G-Data are identical in every wave and count as one source.

\section{Coverage Gap between TLScope and Public Feeds}

Four out of five TLScope-identified malicious servers are not flagged by any of the six providers
within five weeks of their identification. The direction agrees with every prior comparison: MANTIS
finds about five times as many malicious domains as VirusTotal \cite{deniz2025mantis}, the union of
15 public blacklists covered less than 20\% of the malicious domains of most malware families
studied by Kührer et al.\ \cite{kuhrer2014paintitblack}, and Li et al.\ \cite{li2019teaLeaves} and
Bouwman et al.\ \cite{bouwman2020differentcup} found most indicators unique to a single feed. The
magnitudes are not directly comparable --- units (IPs, domains), populations and windows differ ---
but the qualitative picture is consistent: no single feed, and not even the union of many, covers
what a specialised detector finds.

The window-sensitivity result (\S\ref{sec:window_sensitivity}) adds the temporal side of that
picture. Coverage is not fixed: following wave 8 for 60 instead of 39 days lowers its never-detected
rate by 14 percentage points. For an operator, the relevant question is therefore not only whether
feeds will list a server but whether they will list it while the server is still in use --- and
EXPOSURE's measured 1.2-day lifespan for malicious domains \cite{bilge2011exposure} suggests that for
short-lived infrastructure a listing two to three weeks later may come too late to matter.

\section{The Role of the Aggregator}

VirusTotal's engines account for 95.5\% of all detected IPs, and the five standalone providers add
only 4.5\%. In practice, then, the question ``do public feeds know about this server?'' is answered
almost entirely by the engines VirusTotal aggregates --- and within those, by a small and changing
set: twelve engines rotate through the 35 top-five slots, and a single engine (Criminal IP) both
produced the largest bulk listing in the study and changed its retention behaviour mid-study. An
analysis that treated VirusTotal as one provider, or assumed its engines' behaviour to be stable,
would miss both effects.

\section{Indicator Expiry Mechanisms}
\label{sec:expiry}

A flag can disappear because the underlying threat ended, because a provider's retention policy
aged it out, or because a provider's classifier changed its verdict (concept drift
\cite{li2025conceptdrift}). This thesis cannot attribute an individual expiry to one of these causes,
since it observes only the provider's verdict and not the server. The evidence is nonetheless
suggestive at the level of engines: Criminal IP's simultaneous change across four consecutive waves
and AbuseIPDB's lifetime bound match retention policy, while ThreatFox's one-sample burst in wave 3
--- listed and withdrawn within three days --- looks more like a correction than a threat ending.
Distinguishing the causes directly would require observing the servers themselves alongside the
feeds (Chapter~\ref{ch:future_work}).

\section{Threats to Validity}
\label{sec:threats_to_validity}

\subsection{Internal Validity}

\textbf{Ground-truth imperfection.} TLScope reports 91\% precision for malicious server detection,
so up to about 9\% of the tracked IPs could be benign servers that no provider should flag. Such IPs
inflate the never-detected rate for reasons unrelated to feed coverage, so 82.5\% overstates the
coverage gap on genuinely malicious servers by up to about nine percentage points. A benign IP
flagged for an unrelated reason (e.g.\ a shared host with a malicious neighbour) would bias lag and
lifetime in an unknown direction.

\textbf{Right-censoring.} Both lag and lifetime are truncated by the window. Undetected IPs may be
detected later (wave 8 shows that many are, \S\ref{sec:window_sensitivity}), and 64.6\% of measured
lifetimes are still running at the end. All averages in Tables~\ref{tab:top_engine}
and~\ref{tab:provider_coverage} are therefore lower bounds, and comparisons between engines with
different censored shares understate the persistence of the more-censored engine.

\textbf{Sampling granularity.} Lag and lifetime are resolved to the sampling interval: three days in
six waves, two days in wave 5. A ``zero-lag'' IP was flagged within one interval of TLScope's
verdict, not necessarily on the same day, and a flag that appears and disappears between two samples
is invisible.

\textbf{Inconsistent provider coverage.} Censys was queried in waves 2--4 and the start of wave 5
only (\S\ref{sec:provider_availability}). Because Censys flagged just 16 IPs where it was active, its
absence can explain at most a small part of the higher never-detected rate in later waves, but it
does remove one chance of detection per IP from wave 5 onwards.

\textbf{Collection-tool defects.} The audit in \S\ref{sec:data_quality} found one defect that
invalidated part of the pilot wave and led to its exclusion; the remaining waves were checked for
the same defect and are clean. Incomplete VirusTotal responses (0.5--3.0\% of rows) can only
under-count detections.

\subsection{Construct Validity}

\textbf{IP churn.} An IP address is not a stable proxy for a single server: hosting providers reassign
addresses, and a malicious IP can turn benign (or the reverse) within the window. Neither TLScope's
snapshot nor the providers' verdicts can distinguish a change of tenant from a change of
classification.

\textbf{No benign control set.} All tracked IPs come from a malicious-only ground truth. The thesis
measures recall-like coverage but not false-positive rates, since no known-benign IPs are tracked in
parallel.

\textbf{Provider-internal drift.} The set of engines VirusTotal aggregates, and each engine's own
policy, can change during the study --- Criminal IP's mid-study change is one documented example. Such
changes appear as shifts in coverage or lifetime with no change in the underlying threat.

\subsection{External Validity}

\textbf{Detection-methodology selection bias.} TLScope selects on TLS handshake behaviour, so
malicious infrastructure it does not surface (non-TLS abuse, servers mimicking benign TLS stacks) is
outside this study. The coverage gap describes the ecosystem's response to what active TLS
fingerprinting finds --- possibly the very subset that report- and signature-based feeds are least
suited to catch. The study is also IPv4-only.

\textbf{Single vantage point.} All queries are issued from one network location with one API key per
provider. Responses could differ by client identity, subscription tier or query origin.

% ============================================================
\chapter{Future Work}
\label{ch:future_work}
% ============================================================

\begin{description}
  \item[Longer and fixed-length windows] The wave-8 sensitivity check shows coverage still growing
    after five weeks. Following every wave for a fixed, longer period (e.g.\ 90 days) would measure
    eventual coverage and reduce right-censoring of lifetimes.
  \item[Restore or replace Censys] Re-enabling Censys, or replacing it with a comparable
    active-scanning feed, would restore six-provider coverage and remove the confound in
    \S\ref{sec:threats_to_validity}.
  \item[Provider health monitoring and regression tests] Both the Censys deactivation and the
    ThreatFox classification defect were found only by auditing the collected data afterwards.
    Alerting when a configured provider stops returning results, and testing each report processor
    against recorded positive and negative responses, would catch such issues while collection is
    running.
  \item[Re-probing the servers] Re-probing tracked servers (e.g.\ with TLScope
    itself) at each sample would distinguish expiry because the threat ended from expiry by
    retention policy or classifier change (\S\ref{sec:expiry}).
  \item[Benign control set] Tracking a matched set of known-benign IPs would allow false-positive
    rates to be estimated alongside coverage.
  \item[Additional providers] Shodan, AlienVault OTX or commercial feeds would test whether the
    aggregator's dominance holds beyond the providers studied here.
  \item[IPv6 extension] Extending tracking to IPv6 using 6Sense-style scanning
    \cite{williams2024sixsense}.
  \item[Infrastructure-level explanations] AS-level features in the spirit of BEAM
    \cite{chen2024beam}, latency-based clustering in the spirit of CalcuLatency
    \cite{ramesh2024calcuLatency}, or MANTIS-style infrastructure graphs \cite{deniz2025mantis} could
    test whether IPs on the same infrastructure share lag and lifetime --- as the bulk listing events
    suggest --- and predict which IPs a provider is likely to flag.
  \item[Density-aware validation] Applying Tian et al.'s density measurement
    \cite{tian2025density} to provider feature spaces, where accessible, would test whether
    never-detected IPs occupy sparse regions.
  \item[Indicator quality scoring] RAPIER-style \cite{qing2024rapier} confidence scoring, discounting
    correlated engine pairs such as BitDefender/G-Data, to separate high-confidence from
    single-source detections.
\end{description}

% ============================================================
\chapter{Conclusion}
\label{ch:conclusion}
% ============================================================

This thesis set out to measure how public threat intelligence responds to malicious servers that an
independent system has already found. Using dated batches of TLScope-classified malicious IPs as
ground truth, it followed 15\,837 IPs across six providers --- and VirusTotal's 90--100 engines
individually --- every two to three days for five weeks from the day of their classification.

Within that window, 82.5\% of the servers were never flagged by anyone. When servers were flagged,
it was either immediately (47.7\% of detected IPs were already listed when TLScope found them) or
typically one to three weeks later, often in bulk listing events. Most flags, once raised, outlasted
the window, so lifetimes are lower bounds; the broadest engine, Criminal IP, changed mid-study from
keeping its flags to dropping them within one to two weeks. Engines agree rarely --- 83.1\% of
detected IPs rest on a single engine --- yet more often than chance, so the low agreement reflects
how little each engine sees rather than a division of labour. Almost all coverage (95.5\%) comes
from VirusTotal's engines.

Two qualifications matter as much as the numbers. First, coverage depends on how long one waits:
the same wave followed for 60 days instead of 39 shows 68.2\% instead of 82.6\% never detected.
Second, the data behind any such measurement need auditing --- in this study the audit found a
provider deactivated mid-study and a classification defect that led to the exclusion of the pilot
wave. With those caveats stated, the results give a quantitative, ground-truth-anchored account of
where public threat intelligence feeds fall short of specialised active scanning: not only in what
they cover, but in when they cover it and for how long.

% ============================================================
%  BACK MATTER
% ============================================================

\appendix

\chapter{Configuration Reference}
\label{app:config}

API keys are stored in a \texttt{.env} file in the project root. Copy \texttt{.env\_ex} and fill
in the required credentials:
\begin{lstlisting}[language=bash, caption={Example .env configuration}]
VIRUSTOTAL = "VIRUSTOTAL Key::limit"
ABUSEIPDB  = "ABUSEIPDB Key::limit"
CENSYS     = "CENSYS Key||secret::limit"
THREATFOX  = "THREATFOX Key::limit"
# No API key needed for these services; set to "local"
CINSSCORE  = "local::10000"
OPENPHIS   = "local::10000"
\end{lstlisting}

Each value follows the form \texttt{key||secret::rate\_limit}. \texttt{runfetch.py} splits on
\texttt{::} to separate the credential from the per-service rate limit, then on \texttt{||} to
separate key from secret for the one provider that needs both (Censys). Cinsscore and OpenPhish
require no API key --- they use publicly accessible list endpoints --- but still need an entry, set
to \texttt{local}, because a service is registered by the presence of its environment variable.
The OpenPhish variable is spelled \texttt{OPENPHIS} in the codebase; the name must match exactly.
Censys is present in \texttt{.env\_ex} but is not read in the current configuration, since it is
commented out of the active service list (\S\ref{sec:provider_availability}).

\chapter{CLI Reference}
\label{app:cli}

\section{runfetch.py}

\begin{lstlisting}[language=bash]
python runfetch.py -f <folder> [-w WORKERS] [-vw VT_WORKERS]
                   [-q QUEUE_SIZE] [-sf START_FROM] [-dbg]
\end{lstlisting}

\texttt{-w} sets the number of producer workers (Chapter~\ref{ch:implementation}); \texttt{-vw}
sets the number of VirusTotal-polling consumer workers, independently of \texttt{-w}.

\section{runanalysis.py}

\begin{lstlisting}[language=bash]
python runanalysis.py -f <folder> [-t TOP_ENGINES] [-d DAY_DIFF] [-w WORKERS]
                      [-q QUEUE_SIZE] [-o OUTPUT_DIR] [--do_group_graphs]
\end{lstlisting}

\texttt{-t} sets how many of the highest-peak engines get per-engine figures and summary rows
(default 5, the value used for every analysis in this thesis; the per-wave value is recorded as
\texttt{top\_engines} in each wave's \texttt{out\_logs.txt}). It affects only which engines are
singled out: aggregate, agreement and baseline results are computed over all engines.
\texttt{--do\_group\_graphs} additionally plots a per-IP figure for every group; it is off by default
because it produces one figure per group. \texttt{-o} is resolved relative to \texttt{-f}, and the
output directory is emptied before each run.

Configuration overrides are read at two levels. A \texttt{config.json} placed directly inside
\texttt{<folder>} overrides the corresponding command-line value for the whole run, for any key
matching a CLI argument name. A \texttt{config.json} inside an individual \emph{sub-group}
directory overrides that key for that sub-group only; unknown keys are reported and ignored, and a
missing file or key falls back to the CLI/default value. The per-sub-group level is how wave 5's
two-day sampling interval is given to the analysis while the run default stays at three days; the
effective per-sub-group values are echoed into \texttt{out\_logs.txt} as \texttt{group\_overrides},
so the interval used for each sub-group is recoverable from the analysis output alone.

% --- Bibliography ---
\cleardoublepage
\addcontentsline{toc}{chapter}{Bibliography}

\begin{thebibliography}{99}

\bibitem{theofanous2024tlscope}
A.~Theofanous, E.~Papadogiannaki, A.~Shevtsov, and S.~Ioannidis,
``Fingerprinting the Shadows: Unmasking Malicious Servers with Machine Learning-Powered TLS Analysis,''
in \textit{Proc.\ ACM Web Conference (WWW)}, 2024, pp.\ 1933--1944.
DOI:~10.1145/3589334.3645719.

\bibitem{kuhrer2014paintitblack}
M.~Kührer, C.~Rossow, and T.~Holz,
``Paint It Black: Evaluating the Effectiveness of Malware Blacklists,''
in \textit{Proc.\ 17th International Symposium on Research in Attacks, Intrusions and Defenses
(RAID)}, 2014.

\bibitem{li2019teaLeaves}
V.~G.~Li, M.~Dunn, P.~Pearce, D.~McCoy, G.~M.~Voelker, S.~Savage, and K.~Levchenko,
``Reading the Tea Leaves: A Comparative Analysis of Threat Intelligence,''
in \textit{Proc.\ 28th USENIX Security Symposium}, 2019, pp.\ 851--867.

\bibitem{bouwman2020differentcup}
X.~Bouwman, H.~Griffioen, J.~Egbers, C.~Doerr, B.~Klievink, and M.~van~Eeten,
``A different cup of TI? The added value of commercial threat intelligence,''
in \textit{Proc.\ 29th USENIX Security Symposium}, 2020.

\bibitem{zhu2020labeldynamics}
S.~Zhu, J.~Shi, L.~Yang, B.~Qin, Z.~Zhang, L.~Song, and G.~Wang,
``Measuring and Modeling the Label Dynamics of Online Anti-Malware Engines,''
in \textit{Proc.\ 29th USENIX Security Symposium}, 2020.

\bibitem{antonakakis2010notos}
M.~Antonakakis, R.~Perdisci, D.~Dagon, W.~Lee, and N.~Feamster,
``Building a Dynamic Reputation System for DNS,''
in \textit{Proc.\ 19th USENIX Security Symposium}, 2010, pp.\ 273--290.

\bibitem{bilge2011exposure}
L.~Bilge, E.~Kirda, C.~Kruegel, and M.~Balduzzi,
``EXPOSURE: Finding Malicious Domains Using Passive DNS Analysis,''
in \textit{Proc.\ NDSS 2011}, Feb.\ 2011.

\bibitem{deniz2025mantis}
F.~Deniz, M.~Nabeel, T.~Yu, and I.~Khalil,
``MANTIS: Detection of Zero-Day Malicious Domains Leveraging Low Reputed Hosting Infrastructure,''
in \textit{Proc.\ 2025 IEEE Symposium on Security and Privacy (S\&P)}, 2025, pp.\ 1789--1807.
DOI:~10.1109/SP61157.2025.00067.

\bibitem{li2025conceptdrift}
A.~S.~Li, A.~Iyengar, A.~Kundu, and E.~Bertino,
``Revisiting Concept Drift in Windows Malware Detection: Adaptation to Real Drifted Malware with Minimal Samples,''
in \textit{Proc.\ 32nd NDSS 2025}, Feb.\ 2025.

\bibitem{qing2024rapier}
Y.~Qing, Q.~Yin, X.~Deng, Y.~Chen, Z.~Liu, K.~Sun, K.~Xu, J.~Zhang, and Q.~Li,
``Low-Quality Training Data Only? A Robust Framework for Detecting Encrypted Malicious Network Traffic,''
in \textit{Proc.\ 31st NDSS 2024}, Feb.\ 2024.

\bibitem{tian2025density}
J.~Tian, W.~Kong, D.~Gao, T.~Wang, T.~Gu, K.~Qiu, Z.~Wang, and X.~Kuang,
``Density Boosts Everything: A One-stop Strategy for Improving Performance, Robustness, and Sustainability of Malware Detectors,''
in \textit{Proc.\ 32nd NDSS 2025}, Feb.\ 2025.

\bibitem{williams2024sixsense}
G.~Williams, M.~Erdemir, A.~Hsu, S.~Bhat, A.~Bhaskar, F.~Li, and P.~Pearce,
``6Sense: Internet-Wide IPv6 Scanning and its Security Applications,''
in \textit{Proc.\ 33rd USENIX Security Symposium}, 2024, pp.\ 2281--2298.

\bibitem{chen2024beam}
Y.~Chen, Q.~Yin, Q.~Li, Z.~Liu, K.~Xu, Y.~Xu, M.~Xu, Z.~Liu, and J.~Wu,
``Learning with Semantics: Towards a Semantics-Aware Routing Anomaly Detection System,''
in \textit{Proc.\ 33rd USENIX Security Symposium}, 2024, pp.\ 5143--5160.

\bibitem{ramesh2024calcuLatency}
R.~Ramesh, P.~Winter, S.~Korman, and R.~Ensafi,
``CalcuLatency: Leveraging Cross-Layer Network Latency Measurements to Detect Proxy-Enabled Abuse,''
in \textit{Proc.\ 33rd USENIX Security Symposium}, 2024, pp.\ 2263--2280.

\bibitem{zhang2024staleglue}
Y.~Zhang, B.~Liu, H.~Duan, M.~Zhang, X.~Li, F.~Shi, C.~Xu, and E.~Alowaisheq,
``Rethinking the Security Threats of Stale DNS Glue Records,''
in \textit{Proc.\ 33rd USENIX Security Symposium}, 2024, pp.\ 1261--1277.

\bibitem{zhang2024resolverfuzz}
Q.~Zhang, X.~Bai, X.~Li, H.~Duan, Q.~Li, and Z.~Li,
``ResolverFuzz: Automated Discovery of DNS Resolver Vulnerabilities with Query-Response Fuzzing,''
in \textit{Proc.\ 33rd USENIX Security Symposium}, 2024, pp.\ 4729--4746.

\bibitem{jarm2020}
Salesforce Engineering,
``JARM: A Solid TLS Fingerprinting Tool,''
\url{https://engineering.salesforce.com/easily-identify-malicious-servers-on-the-internet-with-jarm-e095edac525a}, 2020.

\end{thebibliography}

\end{document}
